La comprensión de las \emph{construcciones tradicionales} hñähñu como parte de un territorio vivido exige reconocer que su materialidad no se explica solo por lo que se ve. Detrás de un muro de piedra, un entramado vegetal o un recubrimiento de tierra hay decisiones técnicas acumuladas, acuerdos comunitarios, disponibilidades ambientales y memorias de uso. En el Valle del Mezquital, además, muchas de las piezas que permanecen en pie no corresponden a la vivienda completa, sino a \emph{construcciones} asociadas a la vida cotidiana (por ejemplo, cocinas o bodegas) que, al transformarse los modos de habitar y los materiales de mercado, quedan expuestas a sustituciones, reconfiguraciones y lecturas externas que tienden a reducirlas a ``lo antiguo'' o ``lo precario''.

En ese contexto, el presente capítulo reúne los fundamentos conceptuales necesarios para leer el material empírico de esta investigación (fotografías, relatos hablados y cartografía) sin convertirlo en un inventario meramente formal. El objetivo es construir un marco que permita describir y clasificar sistemas constructivos, pero también interpretar sus sentidos y tensiones: qué se conserva, qué cambia, qué se reemplaza y bajo qué condiciones circula el saber hacer.

Para ello, se retoman aportes de la historiografía crítica, los estudios del patrimonio y la teoría de la arquitectura vernácula con cuatro énfasis. Primero, delimitar qué se entiende por \emph{tradición constructiva} y por qué su continuidad no depende únicamente de ``técnicas'' aisladas, sino de cadenas de transmisión situadas. Segundo, situar las construcciones hñähñu dentro del paradigma del \emph{patrimonio vernáculo biocultural}, donde convergen valores materiales, simbólicos y ecológicos. Tercero, precisar tipologías de valor patrimonial y modelos de transmisión--ruptura del saber que ayuden a explicar transformaciones observables (por ejemplo, cambios de cubierta o de acabados) sin suponer automáticamente una pérdida total o una autenticidad fija. Cuarto, considerar las tensiones entre práctica viva, mercado y marcos normativos, identificando categorías analíticas que vinculen escalas locales y globales.

El capítulo se organiza en siete apartados. La Sección~\ref{subsec:tradicion_constructiva} define la tradición constructiva como patrimonio inmaterial encarnado en prácticas y objetos; la Sección~\ref{subsec:vernaculo_biocultural} profundiza en la noción de patrimonio vernáculo biocultural como articulación entre cultura y entorno; la Sección~\ref{subsec:tipologias_valor} revisa tipologías de valor que permiten argumentar por qué ciertas construcciones adquieren relevancia más allá de su utilidad inmediata; la Sección~\ref{subsec:modelos_transmision} presenta modelos para comprender continuidad, aprendizaje y ruptura del saber; la Sección~\ref{subsec:sintesis_operacional} traduce los conceptos a una matriz operativa que conecta preguntas de investigación, dimensiones de análisis y evidencias disponibles; finalmente, la Sección de cierre del capítulo integra estas piezas como guía de lectura para el diagnóstico descriptivo y cartográfico.

Con esta base teórica se busca garantizar coherencia interna entre preguntas, metodología y análisis, evitando reduccionismos (tanto tecnicistas como meramente simbólicos) que limiten la comprensión de las construcciones estudiadas.
